% Generated draft: requires mathematical and source review.
\documentclass[11pt]{article}
\usepackage[margin=25mm]{geometry}
\usepackage{fontspec}
\setmainfont{Latin Modern Roman}
\usepackage{amsmath,amssymb,mathtools,unicode-math}
\setmathfont{Latin Modern Math}
\usepackage{enumitem,needspace,xurl,hyperref}
\setlength{\parindent}{0pt}
\setlength{\parskip}{5pt}
\setlength{\emergencystretch}{3em}
\newcommand{\N}{\mathbb{N}}
\newcommand{\Z}{\mathbb{Z}}
\newcommand{\Q}{\mathbb{Q}}
\newcommand{\R}{\mathbb{R}}
\newcommand{\C}{\mathbb{C}}
\newcommand{\F}{\mathbb{F}}
\newcommand{\A}{\mathbb{A}}
\newcommand{\PP}{\mathbb{P}}
\newcommand{\E}{\mathbb{E}}
\newcommand{\eps}{\varepsilon}
\newcommand{\dd}{\,\mathrm{d}}
\begin{document}
\section*{UnsolvedMath ID: 20003281}
\textbf{Source problem code:} catalogue entry
\subsection{Definitions and mathematical statement}
A. Furman's source conjecture concerns either automorphism actions on finite-volume nilmanifolds or translation actions on \(G/\Lambda\), \(\Lambda\) a lattice. It proposes that a ``large'' acting subgroup \(\Gamma\) has only two kinds of ergodic invariant probability: Haar measure and uniform measures on finite \(\Gamma\)-orbits. The quoted word ``large'' is not defined, so there is no single fully quantified global conjecture to verify. A nilmanifold may be interpreted in the usual sense \(N/\Lambda\) for a nilpotent Lie group and lattice, but the exact source convention and automorphism class are not supplied.

For a specified finite-volume homogeneous space, normalize Haar measure to total mass one. The claimed dichotomy for a fixed action means that every ergodic \(\Gamma\)-invariant probability equals that Haar measure or the uniform probability on one finite orbit. It does not assume Haar ergodicity without a hypothesis ensuring it.

The supplied positive special case is precise: for every integer \(d\geq2\), every ergodic \(\operatorname{SL}_d(\Z)\)-invariant Borel probability on
\[
 \mathbb T^d=\R^d/\Z^d
\]
under its standard linear action is either \(m_{\mathbb T^d}\) or uniform on a finite orbit of rational points. The aid reports a direct content-class/Wiener proof but does not supply it or an authoritative citation. This theorem concerns the full stated integer linear group and is not automatically a theorem for all subgroups deemed large.

The supplied obstruction to one weak interpretation uses
\[
 \Gamma=\operatorname{SL}_2(\Z)\times\operatorname{SL}_2(\Z)
 \curvearrowright\mathbb T^4=\mathbb T^2\times\mathbb T^2,
 \qquad(A,B)(x,y)=(Ax,By).
\]
Its real Zariski closure is the block-diagonal \(\operatorname{SL}_2(\R)\times\operatorname{SL}_2(\R)\), semisimple, Zariski connected, and without compact factors. Nevertheless
\[
 \mu=m_{\mathbb T^2}\otimes\delta_0
\]
is invariant and ergodic, supported on the proper positive-dimensional subtorus \(\mathbb T^2\times\{0\}\). It is neither a finite-orbit measure nor ambient Haar measure. Thus interpreting ``large'' as only those three Zariski-closure conditions is insufficient; this does not refute a different intended definition.

More generally, the aid's intermediate-support test is: a proper closed homogeneous orbit carrying a finite normalized homogeneous measure, preserved ergodically by \(\Gamma\), produces an intermediate measure whenever its support is infinite. Such supports must be excluded by any proposed finite-orbit-or-ambient-Haar hypothesis. The test does not assert that every closed orbit is \(\Gamma\)-invariant or ergodic.

The original references to topological analogues [16,185,148,149,74] remain background. The extraction then begins ``4. Equidistribution'' and a separate setup: \(G\) a Lie group, \(\Gamma\) a lattice, \(U=\{u(t)\}\) one-parameter Ad-unipotent, with \(Ux\) dense in \(G/\Gamma\). Its ensuing question is absent. This is subsection spillover, not a new hypothesis on Furman's acting subgroup. No equidistribution statement is invented from it. The supplied toral result and block obstruction are partial scope tests, not a definition or general solution of ``large.
\subsection{Short English statement}
The undefined word ``large'' prevents a single global formulation. The supplied full \(\operatorname{SL}_d(\Z)\) torus action has the proposed dichotomy, while a block-diagonal example defeats a weaker Zariski-closure interpretation. Proper infinite homogeneous supports provide the stated obstruction.
\subsection{Sources}
\begin{itemize}
\item Local UnsolvedMath registry, permanent ID \(20003281\), exact source code \texttt{catalogue entry}; canonical \texttt{aim-other-notes.json}, \texttt{notes[87]}, AIM domain Other, registry category Miscellaneous, tag \texttt{conjecture}.
\item AIM workshop ``Emerging applications of measure rigidity,'' Conjecture 14 attributed to A. Furman: \url{https://aimath.org/WWN/measrigid/measrigid.pdf}. The supplied verification reports official-PDF inspection with the scare quotes retained. No PDF was fetched here.
\item Source-local labels [16], [185], [148], [149], and [74] concern topological analogues; full entries are not supplied. The original literature field is empty, and no authoritative bibliographic citation for the direct toral theorem or obstruction package is given.
\item Research attempt 1: result \texttt{partial\_result}, problem status at run \texttt{unknown}, rigor label \texttt{proved}, novelty confidence \texttt{low}; classification \texttt{PARTIAL-PROGRESS}, registry status \texttt{partially\_solved}, suggested difficulty L4. Created and updated 14 August 2026. These are supplied assessments requiring review.
\item Attribution: Ulam AI / UnsolvedMath Contributors, CC BY 4.0. The machine-generated research aid is not peer reviewed and requires independent verification before citation.
\end{itemize}
\end{document}
